% 依赖项（必须放在主文档导言区，即 \begin{document} 之前）：
% \usepackage{tabularray}
% \UseTblrLibrary{varwidth}
% \usepackage{enumitem}
\begingroup
% \nopagebreak[4]
\enlargethispage{1cm}
\footnotesize
\centering

\DefTblrTemplate{conthead-text}{default}{（续表）}
\DefTblrTemplate{contfoot-text}{default}{续下页}

\begin{longtblr}[
  caption = {通用商用智能体、MCP工具调用协议与端侧硬件隐私计算智能体},
  label   = {tab:commercial_mcp_edge_agent_security},
]{
  width   = \linewidth,
  colspec = {
    |Q[wd=0.07\linewidth,l,m]
    |Q[wd=0.12\linewidth,l,m]
    |X[1.00,l,m]|
  },
  cells   = {l,m},
  colsep  = 3pt,
  row{1}  = {font=\bfseries},
  rowhead = 1,
  rowsep  = 5pt,
  measure = vbox,
  stretch = -1,
}
\hline
研究方向 & 企业/机构 + 核心学者 & 自研项目、产学合作、商用安全产品 \\
\hline

\SetCell[r=3]{l,m} 通用商用智能体安全评估与治理规范
& OpenAI：Sam Altman、Greg Brockman
& \begin{itemize}[leftmargin=*,topsep=0pt,partopsep=0pt,itemsep=1pt,parsep=0pt]
\item \textbf{ChatGPT Agent交互智能体：}搭建外部多专家红队测试网络，Red Teaming Network由16位具备生物安全背景的博士研究员组成，共提交110次攻击尝试，其中16次超过内部风险阈值~\cite{openaiChatGPTAgentSystemCard2025,openaiRedTeamNetwork2025}。
\item \textbf{开源安全护栏体系：}OpenAI Agents SDK内置三层输入/输出/工具护栏，支持检测风险时立即中止运行的tripwire机制，配合运行时Tracing审计与高风险动作的human-in-the-loop审批机制~\cite{openaiAgentsSDK2025,openaiGuardrailsDoc2026,openaiTracingDoc2026,openaiGuardrailsApprovalsDoc2026}。
\item \textbf{全球统一风险披露体系：}System Card持续迭代智能体能力、风险与透明度评估指标，全产品线强制落地，并以公开白皮书系统阐述外部专家参与风险评估的方法论~\cite{openaiChatGPTAgentSystemCard2025,openaiExternalRedTeaming2025}。
\end{itemize} \\
\cline{2-3}
& 英国AI安全研究所UK AISI、FAR.AI独立安全实验室
& \begin{itemize}[leftmargin=*,topsep=0pt,partopsep=0pt,itemsep=1pt,parsep=0pt]
\item \textbf{深度访问链路漏洞审计：}UK AISI获得对ChatGPT Agent内部推理链和策略文本的深度访问权限，完成官方漏洞审计测试，测出7个通用漏洞~\cite{venturebeatChatGPTAgent2025}。
\item \textbf{商用智能体安全缺陷批评：}FAR.AI经40小时独立测试发现3个部分漏洞，提出商用智能体过度依赖运行时监控的结构性安全缺陷~\cite{venturebeatChatGPTAgent2025}。
\end{itemize} \\
\cline{2-3}
& 斯坦福CRFM：Percy Liang
& \begin{itemize}[leftmargin=*,topsep=0pt,partopsep=0pt,itemsep=1pt,parsep=0pt]
\item \textbf{全链路可比较评测理论对齐：}HELM全链路可比较评测理论与OpenAI Agents SDK的Tracing架构在方法论上形成理论对应，共同强调把系统行为拆解为可追溯、可比较的具体环节~\cite{openaiTracingDoc2026}。
\end{itemize} \\
\hline

\SetCell[r=3]{l,m} 智能体工具调用MCP协议与目标对齐安全
& Anthropic：Dario Amodei、Jared Kaplan
& \begin{itemize}[leftmargin=*,topsep=0pt,partopsep=0pt,itemsep=1pt,parsep=0pt]
\item \textbf{Agent全生态工具交互协议：}MCP工具通信协议基于JSON-RPC 2.0构建，支撑Claude Agent全生态工具交互，截至2026年初已形成超1万活跃服务器、17.7万注册工具、月均9700万次SDK下载的生态规模~\cite{mcpFormalFramework2026}。
\item \textbf{高危分级智能体系统：}Claude Code代码智能体、Mythos高危分级智能体系统均采用分层行为约束护栏~\cite{anthropicClaudeCodeSecurity2026}。
\item \textbf{协议架构安全责任取舍：}面对OX Security披露的架构级RCE漏洞，Anthropic确认该行为符合设计预期，以维持协议灵活性为由拒绝修复，将输入净化责任划给下游开发者~\cite{oxSecurityMCP2026}。
\end{itemize} \\
\cline{2-3}
& 加州伯克利CHAI：Stuart Russell
& \begin{itemize}[leftmargin=*,topsep=0pt,partopsep=0pt,itemsep=1pt,parsep=0pt]
\item \textbf{人类兼容AI理论与目标对齐：}Russell作为Berkeley CHAI创始主任，其人类兼容AI理论关注目标不确定性和系统可控性，与高自治Agent的目标漂移、价值对齐问题存在直接理论关联~\cite{berkeleyCHAI2026}。
\end{itemize} \\
\cline{2-3}
& OX Security、Check Point安全厂商；密码学形式化验证学界
& \begin{itemize}[leftmargin=*,topsep=0pt,partopsep=0pt,itemsep=1pt,parsep=0pt]
\item \textbf{MCP协议远程代码执行漏洞披露：}OX Security披露MCP协议"设计如此"的远程代码执行漏洞，影响超1.5亿次下载和20万台服务器，攻陷11个主流MCP市场中的9个，产生至少10个高危或严重级别CVE编号~\cite{oxSecurityMCP2026,thehackernewsMCP2026}。
\item \textbf{Claude Code代码执行漏洞修复：}Check Point于2025年9至10月披露Claude Code三项漏洞（CVE-2025-59536，CVSS 8.7；CVE-2026-21852，CVSS 5.3），由厂商协同完整修复，分别在1.0.111版本和2.0.65版本发布安全补丁~\cite{checkpointClaudeCode2026,githubClaudeCodeAdvisory2026}。
\item \textbf{MCP威胁分类与形式化验证：}学术界已提出MCP的威胁分类模型，并指出借鉴Tamarin、ProVerif等密码协议验证工具对MCP做形式化验证仍是开放且困难的研究方向~\cite{mcpFormalFramework2026,mcpSoK2026}。
\end{itemize} \\
\hline

\SetCell[r=3]{l,m} 端侧与硬件机密计算隐私智能体安全技术
& Apple：Craig Federighi 隐私安全实验室
& \begin{itemize}[leftmargin=*,topsep=0pt,partopsep=0pt,itemsep=1pt,parsep=0pt]
\item \textbf{端云分层隐私架构：}Apple Intelligence移动端本地智能体采用端云分层数据最小化架构~\cite{applePCCSecurityGuide2026}。
\item \textbf{可验证云端推理平台：}Private Cloud Compute（PCC）通过Virtual Research Environment向全球研究者开放本地复现能力，配合开源形式化验证工具、函数库与安全悬赏计划，构建透明可信推理体系~\cite{appleSecurityResearch2026}。
\end{itemize} \\
\cline{2-3}
& EPFL SPRING Lab：Carmela Troncoso
& \begin{itemize}[leftmargin=*,topsep=0pt,partopsep=0pt,itemsep=1pt,parsep=0pt]
\item \textbf{隐私工程理论支撑端侧架构：}Troncoso团队长期研究隐私工程、数据最小化和匿名化技术，相关理论与端侧/云端隐私智能体架构问题高度相关~\cite{springLab2026}。
\end{itemize} \\
\cline{2-3}
& NVIDIA：Jensen Huang，NYU；Ramesh Karri
& \begin{itemize}[leftmargin=*,topsep=0pt,partopsep=0pt,itemsep=1pt,parsep=0pt]
\item \textbf{GPU硬件机密计算平台：}Confidential Agent GPU硬件机密计算平台自H100起首次支持机密计算，通过片上硬件信任根实现安全启动和远程认证，在Blackwell架构上进一步延伸~\cite{nvidiaConfidentialComputingProduct2026}。
\item \textbf{硬件侧信道安全研究支撑：}Karri教授长期研究硬件木马、芯片供应链安全和侧信道分析，与GPU侧信道信息泄露防护议题直接相关，学术界也已对TEE模式下大模型推理性能开销做独立评测~\cite{karriNYU2026}。
\end{itemize} \\
\hline
\end{longtblr}
\endgroup
